
\iffalse

Some scenarios additionally benefit from splitting the multiplication of sparse matrices and other routines operating on sparse data structures into a symbolic and a numeric stage. This, for example, can allow for additional optimization in the repeated multiplication of matrices, changing only the numeric values but preserving the sparsity pattern. Hence, we propose a variant of the above interface, splitting the routines
\texttt{multiply\_compute} and
\texttt{multiply\_fill} into two:

\begin{minted}{c++}
using namespace spblas;
// csr_wrapper<float> A,B,D filled elsewhere
// csr_wrapper<float> C(c_nrows, c_ncols);

operation_state state(allocator);
multiply_inspect(policy, state, A, B, C, D); // optional
multiply_symbolic_compute(policy, state, A, B, C, D);
index_t nnz = state.get_result_nnz();
// allocate C arrays and put in C (possibly both structures
// and values or just structure arrays at this point)
multiply_symbolic_fill(policy, state, A, B, C, D);
// C structure is ready to be used

// allocate, if necessary, value arrays and place them in C
multiply_numeric_compute(policy, state, A, B, C, D);
multiply_numeric_fill(policy, state, A, B, C, D);
// C structure and values are now able to be used
\end{minted}



Because we do not specify at which phase the output matrix is prepared internally, the \textit{compute phase} (e.g. \texttt{multiply\_numeric\_compute}) and the \textit{fill phase} (e.g. \texttt{sparse\_multiply\_numeric\_fill}) should be regarded as a pair and for repeated execution, both phases should be called.  It is up to implementations to detect on repeated calls which internal data in the \texttt{multiply\_state\_t} object is reusable and what must be recomputed.

We want to provide some examples for the use of this API. 
\begin{enumerate}
    \item Single use:
\begin{minted}{c++}

operation_state state(allocator);
multiply_inspect(policy, state, A, B, C, D); // optional
multiply_compute(policy, state, A, B, C, D);
index_t nnz = state.get_result_nnz();
// allocate C arrays 
multiply_fill(policy, state, A, B, C, D);
// C structure is now able to be used
\end{minted}
\item 
Repeated invocation of the functionality with numerically differing matrices:
\begin{minted}{c++}
operation_state state(allocator);
multiply_inspect(policy, state, A, B, C, D); // optional
for (i=0 i<n; i++) {
  // may do nothing if nnz is already computed
  // and information is stored in the state
  multiply_compute(policy, state, A, B, C, D);
  index_t nnz = state.get_result_nnz();
  if (i == 0 }{
    // allocate C arrays 
  }
  multiply_fill(policy, state, A, B, C, D);
  // C structure is now able to be used
}
\end{minted}
\item 
Using the split into numeric and symbolic phases for repeated invocation of the functionality with numerically differing matrices:
\begin{minted}{c++}

operation_state state(allocator);
multiply_inspect(policy, state, A, B, C, D); // optional
multiply_symbolic_compute(policy, state, A, B, C, D);
index_t nnz = state.get_result_nnz();
// allocate C arrays 
multiply_symbolic_fill(policy, state, A, B, C, D);
for (i=0 i<n; i++) {
  // only numeric computations inside the loop
  multiply_numeric_compute(policy, state, A, B, C, D);
  multiply_numeric_fill(policy, state, A, B, C, D);
  // C structure is now able to be used
}
\end{minted}
\end{enumerate}

\begin{figure}
    \centering
    \includegraphics[width=.9\columnwidth]{graphics/different_SpGEMM.png}
    \caption{Different workflows for the CSR-based SpGEMM.}
    \label{fig:SpGEMM_workflows}
\end{figure}

\fi
